%!TEX root = ../thesis.tex
% ******************************* Thesis Appendix E ********************************
\chapter{Source Data for the Swept Figures}
\label{app:sourcedata}

\noindent Appendix~\ref{app:hyperparameters} gives the \emph{inputs} of each experiment
and Appendix~\ref{app:code} the code that runs them. This appendix gives the
\emph{outputs}: the per-cell numerical values plotted in every swept figure of
Chapters~\ref{ch:model},~\ref{ch:learning} and~\ref{ch:conclusion}, so that each figure
can be \emph{redrawn} without re-running the simulation, and so that the exact means and
confidence intervals behind the curves are on record. Each table is the data of one
figure; the caption names the figure and the reference run. Unless noted, every value
is a mean over $20$ random seeds with a percentile-bootstrap $95\%$ confidence interval
$[\text{lo},\text{hi}]$ ($10^4$ resamples; Listing~\ref{lst:siox-src-stats}), the ensemble and
interval convention used throughout the results chapters. The array-scale feasibility
grids (Tables~\ref{table:E7},~\ref{table:E8} and~\ref{table:E9}) use $12$ seeds, as
stated in their captions. These are the values written by the reference code; minor variation across
BLAS/NumPy versions is expected in the last reported digit.

% =====================================================================
\section{Retention sets the credit window}

\begin{table}[H]
    \caption[Source data: trace-level credit window (Fig.~\ref{fig:7c}).]{Learned
    cue-to-distractor weight ratio against action--reward delay $D$ (s), for three
    retention times and the no-trace control; the data of Figure~\ref{fig:7c}. Values
    are mean\,[lo,\,hi] over $20$ seeds. A ratio above the learned threshold ($2.0$)
    marks a bridged delay.}
    \begin{center}
    \footnotesize
    \begin{tabularx}{\textwidth}{@{}l *{8}{Y}@{}}
    \toprule
    $D$ (s) & 1 & 2 & 5 & 10 & 20 & 40 & 70 & 100 \\
    \midrule
    $\tau_{leak}{=}20$ & 1.26 & 2.00 & 4.24 & 4.25 & 4.07 & 3.83 & 3.66 & 3.66 \\
    \;[lo,\,hi] & \tiny[1.19,1.33] & \tiny[1.89,2.12] & \tiny[3.96,4.54] & \tiny[3.90,4.61] & \tiny[3.68,4.47] & \tiny[3.50,4.18] & \tiny[3.30,4.05] & \tiny[3.33,4.01] \\
    $\tau_{leak}{=}5$ & 2.99 & 4.27 & 4.19 & 3.93 & 3.82 & 2.38 & 0.99 & 1.01 \\
    \;[lo,\,hi] & \tiny[2.76,3.24] & \tiny[3.93,4.62] & \tiny[3.85,4.56] & \tiny[3.65,4.25] & \tiny[3.60,4.03] & \tiny[2.27,2.49] & \tiny[0.94,1.04] & \tiny[0.95,1.07] \\
    $\tau_{leak}{=}1$ & 3.99 & 3.95 & 3.86 & 1.20 & 1.02 & 0.99 & 1.01 & 1.00 \\
    \;[lo,\,hi] & \tiny[3.85,4.15] & \tiny[3.79,4.12] & \tiny[3.76,3.96] & \tiny[1.17,1.22] & \tiny[1.00,1.05] & \tiny[0.96,1.03] & \tiny[0.97,1.04] & \tiny[0.97,1.03] \\
    no-trace & 1.00 & 1.01 & 1.00 & 1.00 & 1.00 & 1.00 & 1.00 & 1.00 \\
    \bottomrule
    \end{tabularx}
    \label{table:E1}
    \end{center}
    \vspace*{-\baselineskip}
\end{table}

\begin{table}[H]
    \caption[Source data: spiking cue saturation (Fig.~\ref{fig:7d}).]{Saturation of
    the cue synapse toward its conductance bound against delay $D$ (s), in the full
    spiking network, for three retention times; the data of Figure~\ref{fig:7d}. Mean
    over $20$ seeds (values near zero indicate no credit assigned).}
    \begin{center}
    \footnotesize
    \begin{tabularx}{\textwidth}{@{}l *{6}{Y}@{}}
    \toprule
    $D$ (s) & 1 & 2 & 5 & 10 & 20 & 40 \\
    \midrule
    $\tau_{leak}{=}10$ & 0.084 & 0.237 & 0.804 & 0.974 & 0.924 & $-0.005$ \\
    $\tau_{leak}{=}2$  & 0.023 & 0.034 & 0.030 & 0.006 & $-0.001$ & $-0.001$ \\
    $\tau_{leak}{=}0.5$ & 0.001 & 0.000 & $-0.000$ & $-0.000$ & $-0.000$ & $-0.000$ \\
    \bottomrule
    \end{tabularx}
    \label{table:E2}
    \end{center}
    \vspace*{-\baselineskip}
\end{table}

\begin{table}[H]
    \caption[Source data: bandit reward rate, delay $\times$ retention (Fig.~\ref{fig:7e}).]{Final
    reward rate of the two-state contextual bandit against delay $D$ (s) for three
    retention times; the delay--retention data underlying Figure~\ref{fig:7e}. Mean\,[lo,\,hi]
    over $20$ seeds. Chance is $0.5$.}
    \begin{center}
    \footnotesize
    \begin{tabularx}{\textwidth}{@{}l *{5}{Y}@{}}
    \toprule
    $D$ (s) & 1 & 2 & 5 & 10 & 20 \\
    \midrule
    $\tau_{leak}{=}10$ & 0.730 & 0.896 & 0.989 & 0.999 & 1.000 \\
    \;[lo,\,hi] & \tiny[0.711,0.749] & \tiny[0.882,0.910] & \tiny[0.984,0.993] & \tiny[0.997,1.000] & \tiny[1.000,1.000] \\
    $\tau_{leak}{=}2$ & 0.634 & 0.691 & 0.726 & 0.568 & 0.512 \\
    \;[lo,\,hi] & \tiny[0.616,0.653] & \tiny[0.670,0.711] & \tiny[0.702,0.752] & \tiny[0.551,0.585] & \tiny[0.492,0.531] \\
    $\tau_{leak}{=}0.5$ & 0.510 & 0.501 & 0.512 & 0.504 & 0.510 \\
    \;[lo,\,hi] & \tiny[0.491,0.529] & \tiny[0.482,0.521] & \tiny[0.487,0.537] & \tiny[0.488,0.520] & \tiny[0.488,0.530] \\
    \bottomrule
    \end{tabularx}
    \label{table:E3}
    \end{center}
    \vspace*{-\baselineskip}
\end{table}

\begin{table}[H]
    \caption[Source data: retention--delay scaling law (Fig.~\ref{fig:7l}).]{Maximum
    learnable delay $D_{\max}$ (s, interpolated at the $0.75$ criterion) against
    retention $\tau_{leak}$ (s), from the dense sequential sweep of Figure~\ref{fig:7l}
    ($20$ seeds). A line through the origin over all points gives $s_D\approx15.4$
    ($R^2=0.99$); no point is censored on the extended delay grid.}
    \begin{center}
    \footnotesize
    \begin{tabularx}{\textwidth}{@{}l *{7}{Y}@{}}
    \toprule
    $\tau_{leak}$ & 1.0 & 1.5 & 2.0 & 2.5 & 3.0 & 4.0 & 5.0 \\
    $D_{\max}$    & 7.3 & 13.8 & 17.5 & 25.3 & 30.2 & 42.2 & 58.6 \\
    \midrule
    $\tau_{leak}$ & 6.5 & 8.0 & 12.0 & 20.0 & 30.0 & 40.0 & \\
    $D_{\max}$    & 82.7 & 102.2 & 160.8 & 287.5 & 478.6 & 637.9 & \\
    \bottomrule
    \end{tabularx}
    \label{table:E4}
    \end{center}
    \vspace*{-\baselineskip}
\end{table}

% =====================================================================
\section{Scaling, sequential credit, and the algorithmic benchmark}

\begin{table}[H]
    \caption[Source data: closed-loop scaling and remedies (Fig.~\ref{fig:7g}).]{(Top)
    Final reward rate against array size $S\times A$ for the device trace and the
    no-trace control, the scaling data of Figure~\ref{fig:7g}a; the per-size
    convergence criterion is $\tfrac{1}{2}(1+1/A)$. (Bottom) The five remedies on the
    failing $8\times8$ case, Figure~\ref{fig:7g}b. Mean\,[lo,\,hi] over $20$ seeds.}
    \begin{center}
    \footnotesize
    \begin{tabularx}{\textwidth}{@{}l Y Y Y c@{}}
    \toprule
    $S\times A$ & Device & no-trace & criterion & pass \\
    \midrule
    $2\times2$  & 0.977\,\tiny[0.970,0.983] & 0.490\,\tiny[0.470,0.512] & 0.750 & yes \\
    $4\times2$  & 0.921\,\tiny[0.909,0.933] & 0.491\,\tiny[0.474,0.510] & 0.750 & yes \\
    $4\times4$  & 0.821\,\tiny[0.803,0.838] & 0.237\,\tiny[0.221,0.254] & 0.625 & yes \\
    $8\times4$  & 0.542\,\tiny[0.526,0.558] & 0.251\,\tiny[0.232,0.271] & 0.625 & no \\
    $8\times8$  & 0.308\,\tiny[0.287,0.331] & 0.127\,\tiny[0.113,0.142] & 0.562 & no \\
    $12\times8$ & 0.219\,\tiny[0.200,0.237] & 0.126\,\tiny[0.112,0.140] & 0.562 & no \\
    \midrule
    \multicolumn{5}{@{}l}{\emph{Remedies on the $8\times8$ case (criterion $0.562$):}} \\
    R0 baseline (device, $1500$ trials)          & 0.308\,\tiny[0.287,0.331] & & & no \\
    R1 larger budget (device, $6000$)            & 0.701\,\tiny[0.664,0.735] & & & yes \\
    R2 directed exploration ($\sigma\,0.45{\to}0.05$) & 0.468\,\tiny[0.449,0.486] & & & no \\
    R3 abstract trace (not device, $1500$)       & 0.751\,\tiny[0.699,0.801] & & & yes \\
    R4 directed $+$ budget ($6000$)              & 0.966\,\tiny[0.951,0.980] & & & yes \\
    \bottomrule
    \end{tabularx}
    \label{table:E5}
    \end{center}
    \vspace*{-\baselineskip}
\end{table}

\begin{table}[H]
    \caption[Source data: T-maze retention $\times$ delay and benchmark (Table~\ref{table:7a}).]{(Top)
    T-maze reward rate against goal-reward delay $D$ (s) for three retention times, the
    retention--delay grid behind the sequential result and $D_{\max}$ points
    (Section~\ref{sec:7p2p3}). (Bottom) The algorithmic benchmark of
    Table~\ref{table:7a}: final reward rate at the ceiling operating point. Mean\,[lo,\,hi]
    over $20$ seeds; chance $0.5$.}
    \begin{center}
    \footnotesize
    \begin{tabularx}{\textwidth}{@{}l *{6}{Y}@{}}
    \toprule
    $D$ (s) & 2 & 4 & 8 & 16 & 32 & 64 \\
    \midrule
    $\tau_{leak}{=}2$  & 0.933 & 0.897 & 0.999 & 0.935 & 0.505 & 0.505 \\
    $\tau_{leak}{=}5$  & 0.914 & 0.988 & 0.988 & 0.975 & 0.995 & 0.616 \\
    $\tau_{leak}{=}20$ & 0.971 & 0.998 & 0.996 & 0.986 & 1.000 & 0.997 \\
    \midrule
    \multicolumn{7}{@{}l}{\emph{Benchmark at the ceiling operating point (Table~\ref{table:7a}):}} \\
    \multicolumn{7}{@{}l}{\quad Device trace $1.000\,[1.000,1.000]$; \ R-STDP $0.997\,[0.994,0.999]$;} \\
    \multicolumn{7}{@{}l}{\quad e-prop $1.000\,[1.000,1.000]$; \ no-trace $\approx$ chance.} \\
    \bottomrule
    \end{tabularx}
    \label{table:E6}
    \end{center}
    \vspace*{-\baselineskip}
\end{table}

% =====================================================================
\section{Deep all-local network and biological grounding}

\begin{table}[H]
    \caption[Source data: deep all-local conditions (Figs.~\ref{fig:7m},~\ref{fig:7dms}).]{Final
    reward rate per condition for the deep all-local XOR network (Figure~\ref{fig:7m},
    $H=32$) and the temporal-distractor convergence task (Figure~\ref{fig:7dms}). Mean\,[lo,\,hi]
    over $20$ seeds; ``seeds solved'' counts seeds reaching the $0.75$ criterion. Chance
    $0.5$.}
    \begin{center}
    \footnotesize
    \begin{tabularx}{\textwidth}{@{}l Y Y c@{}}
    \toprule
    Condition & Reward rate & [lo,\,hi] & seeds solved \\
    \midrule
    \multicolumn{4}{@{}l}{\emph{Deep XOR (Fig.~\ref{fig:7m}):}} \\
    shallow (one trained layer) & 0.315 & \tiny[0.280,0.355] & --- \\
    ELM (fixed random hidden)   & 0.496 & \tiny[0.398,0.608] & --- \\
    global scalar               & 0.504 & \tiny[0.489,0.519] & --- \\
    DFA alone                   & 0.734 & \tiny[0.604,0.855] & 11/20 \\
    DFA $+$ homeostasis         & 1.000 & \tiny[1.000,1.000] & 20/20 \\
    no-trace                    & 0.539 & \tiny[0.472,0.613] & --- \\
    \midrule
    \multicolumn{4}{@{}l}{\emph{Deep XOR $+$ temporal distractor (Fig.~\ref{fig:7dms}):}} \\
    full stack, no distractor   & 0.988 & \tiny[0.965,1.000] & 20/20 \\
    full stack $+$ distractor   & 1.000 & \tiny[1.000,1.000] & 20/20 \\
    no-homeostasis $+$ distractor & 0.577 & \tiny[0.466,0.690] & 4/20 \\
    no-trace $+$ distractor     & 0.539 & \tiny[0.472,0.613] & 4/20 \\
    \bottomrule
    \end{tabularx}
    \label{table:E6b}
    \end{center}
    \vspace*{-\baselineskip}
\end{table}

\begin{table}[H]
    \caption[Source data: Frank task and dopamine capstone (Figs.~\ref{fig:7q},~\ref{fig:7w}).]{Final
    performance on the Frank probabilistic-selection task (Figure~\ref{fig:7q}) and the
    measured-reward capstone (Figure~\ref{fig:7w}). Mean\,[lo,\,hi] over $20$ seeds.
    The EEG-driven and dopamine-driven rows use rewards decoded from external
    recordings (ds003474; \abbrev{DANDI}~000351).}
    \begin{center}
    \footnotesize
    \begin{tabularx}{\textwidth}{@{}X Y Y@{}}
    \toprule
    Condition & Reward rate & [lo,\,hi] \\
    \midrule
    Frank, shallow (probability of higher-value choice) & 0.862 & \tiny[0.837,0.883] \\
    Frank, deep (DFA)                                   & 0.740 & \tiny[0.704,0.778] \\
    Frank, deep (DFA $+$ homeostasis)                   & 0.798 & \tiny[0.741,0.850] \\
    Frank, deep, real-EEG reward                        & 0.639 & \tiny[0.599,0.686] \\
    Frank, deep, shuffled-reward control                & 0.513 & \tiny[0.481,0.543] \\
    Capstone, shallow, synthetic reward                 & 0.84  & \tiny(15/20 conv.) \\
    Capstone, deep (DFA $+$ homeostasis), dopamine reward & 0.92 & \tiny(disjoint from DFA-alone $0.65$) \\
    \bottomrule
    \end{tabularx}
    \label{table:E6c}
    \end{center}
    \vspace{0.2em}
    \noindent\footnotesize The dopamine reward decoder achieves a balanced accuracy of
    $0.885$ across $10$ mouse photometry sessions (nine of ten above $0.69$); the human
    EEG decoder $0.688\pm0.103$ over $5$ subjects.
\end{table}

% =====================================================================
\section{Array-scale feasibility grids (12 seeds)}

\begin{table}[H]
    \caption[Source data: array-scale fault tolerance (Fig.~\ref{fig:7as}).]{Final
    reward rate of the deep all-local rule against hidden width $H$ and stuck-device
    fraction $p$, under the measured \abbrev{SiOx} fault prior (stuck-open $+$
    asymmetric device-to-device spread, no Poole--Frenkel read nonlinearity); the data
    of Figure~\ref{fig:7as}. Mean\,[lo,\,hi] over $12$ seeds. The eligibility-zeroed
    control (no-trace) stays at chance ($\approx0.50$) at every cell.}
    \begin{center}
    \footnotesize
    \begin{tabularx}{\textwidth}{@{}l *{4}{Y}@{}}
    \toprule
    & $p{=}0$ & $p{=}0.05$ & $p{=}0.20$ & $p{=}0.50$ \\
    \midrule
    $H{=}8$   & 0.755\,\tiny[0.737,0.771] & 0.738\,\tiny[0.720,0.756] & 0.712\,\tiny[0.688,0.737] & 0.589\,\tiny[0.563,0.615] \\
    $H{=}32$  & 0.906\,\tiny[0.887,0.926] & 0.888\,\tiny[0.864,0.912] & 0.826\,\tiny[0.806,0.849] & 0.678\,\tiny[0.660,0.695] \\
    $H{=}128$ & 0.964\,\tiny[0.952,0.976] & 0.953\,\tiny[0.943,0.963] & 0.873\,\tiny[0.854,0.893] & 0.737\,\tiny[0.716,0.757] \\
    $H{=}512$ & 0.977\,\tiny[0.971,0.982] & 0.975\,\tiny[0.967,0.983] & 0.905\,\tiny[0.880,0.927] & 0.740\,\tiny[0.721,0.756] \\
    \midrule
    control   & \multicolumn{4}{c}{no-trace $\approx0.49$--$0.51$ across all $H,p$} \\
    \bottomrule
    \end{tabularx}
    \label{table:E7}
    \end{center}
    \vspace*{-\baselineskip}
\end{table}

\begin{table}[H]
    \caption[Source data: hybrid stack at array scale (Fig.~\ref{fig:7hs}).]{Final
    reward rate of the full local stack behind a real frozen perception front-end
    (four-way task, $A=4$, chance $0.25$), against hidden width $H$ and stuck fraction
    $p$; the data of Figure~\ref{fig:7hs}. Mean\,[lo,\,hi] over $12$ seeds, disjoint
    throughout from the eligibility-zeroed control.}
    \begin{center}
    \footnotesize
    \begin{tabularx}{\textwidth}{@{}l Y Y Y@{}}
    \toprule
    & $p{=}0$ & $p{=}0.20$ & $p{=}0.50$ \\
    \midrule
    $H{=}32$  full stack & 0.475\,\tiny[0.463,0.487] & 0.501\,\tiny[0.488,0.514] & 0.503\,\tiny[0.490,0.515] \\
    \quad\ no-trace      & 0.253\,\tiny[0.233,0.273] & 0.250\,\tiny[0.225,0.275] & 0.242\,\tiny[0.222,0.264] \\
    $H{=}128$ full stack & 0.512\,\tiny[0.500,0.525] & 0.526\,\tiny[0.519,0.534] & 0.555\,\tiny[0.546,0.562] \\
    \quad\ no-trace      & 0.246\,\tiny[0.232,0.261] & 0.238\,\tiny[0.218,0.258] & 0.260\,\tiny[0.238,0.282] \\
    $H{=}512$ full stack & 0.506\,\tiny[0.495,0.517] & 0.533\,\tiny[0.527,0.539] & 0.564\,\tiny[0.557,0.570] \\
    \quad\ no-trace      & 0.268\,\tiny[0.247,0.287] & 0.259\,\tiny[0.239,0.276] & 0.255\,\tiny[0.237,0.275] \\
    \bottomrule
    \end{tabularx}
    \label{table:E8}
    \end{center}
    \vspace*{-\baselineskip}
\end{table}

\begin{table}[H]
    \caption[Source data: retention-as-asset distal-cue control (Fig.~\ref{fig:7dc}).]{Final
    reward rate of the deep all-local agent on the single-deposit distal-cue task
    against cue--reward gap (s), for a long-retention trace ($\tau_{leak}=20$), a
    short-retention control ($\tau_{leak}=1.5$), and no-trace; the data of
    Figure~\ref{fig:7dc}. Mean\,[lo,\,hi] over $20$ seeds at the fault-free ($p=0$)
    operating point; chance $0.5$.}
    \begin{center}
    \footnotesize
    \begin{tabularx}{\textwidth}{@{}l *{4}{Y}@{}}
    \toprule
    gap (s) & 1 & 4 & 10 & 20 \\
    \midrule
    long ($\tau_{leak}{=}20$)  & 0.809\,\tiny[0.692,0.909] & 0.912\,\tiny[0.850,0.966] & 0.902\,\tiny[0.839,0.959] & 0.856\,\tiny[0.785,0.923] \\
    short ($\tau_{leak}{=}1.5$) & 0.451\,\tiny[0.400,0.507] & 0.540\,\tiny[0.462,0.626] & 0.497\,\tiny[0.428,0.573] & 0.473\,\tiny[0.421,0.523] \\
    no-trace                   & 0.470\,\tiny[0.431,0.514] & 0.470\,\tiny[0.430,0.513] & 0.470\,\tiny[0.430,0.513] & 0.470\,\tiny[0.430,0.513] \\
    \bottomrule
    \end{tabularx}
    \label{table:E9}
    \end{center}
    \vspace*{-\baselineskip}
\end{table}

% =====================================================================
\section{Exploratory extensions (Chapter~\ref{ch:conclusion})}

\begin{table}[H]
    \caption[Source data: Chapter~\ref{ch:conclusion} extensions
    (Figs.~\ref{fig:8a}--\ref{fig:8h}).]{Per-figure values for the exploratory
    extensions of Chapter~\ref{ch:conclusion}. Mean\,[lo,\,hi] over $20$ seeds
    ($12$ for the longer-horizon grid, Fig.~\ref{fig:8a}); chance $0.5$ unless noted.}
    \begin{center}
    \footnotesize
    \begin{tabularx}{\textwidth}{@{}l X@{}}
    \toprule
    Figure & Values \\
    \midrule
    \ref{fig:8a} long-horizon ($L$) & device $0.99$ across $L{=}3$--$12$ (both 2- and 3-action); e-prop $0.84{\to}0.81$ (2-action), $\approx0.59$ (3-action); no-trace chance. \\
    \addlinespace
    \ref{fig:8b} reversal & pre-reversal: device $0.995$, abstract $1.000$, no-trace $0.500$; post-reversal: device $0.988$\,\tiny[0.986,0.991], abstract $1.000$, no-trace $0.491$. Re-acquisition cost: device $1568$ vs abstract $120$ trials ($13.1\times$). \\
    \addlinespace
    \ref{fig:8c} multi-timescale & hetero-raw $0.537$\,\tiny[0.537,0.555]; hetero-homeo $1.000$\,\tiny[1.000,1.000]; $\tau$-oracle $0.992$; best single $0.771$; no-trace $0.494$. \\
    \addlinespace
    \ref{fig:8d} vector timer & vector $0.768$\,\tiny[0.758,0.778]; scalar $0.503$\,\tiny[0.492,0.515]; no-trace $0.504$\,\tiny[0.493,0.515] (aliasing pair $t_A{=}10.5$, $t_B{=}30.5$\,s). \\
    \addlinespace
    \ref{fig:8g} WM/STC (per $\tau_{leak}$) & WM (cue hold): $0.86$ at $0.8$\,s, $1.00$ for $\tau_{leak}\ge1.3$\,s. STC (credit tag): $0.80$ at $0.8$\,s, $0.82$ at $1.3$\,s, $0.96$ at $3.6$\,s, $1.00$ for $\tau_{leak}\ge6$\,s. \\
    \addlinespace
    \ref{fig:8h} device-TD (by $L$) & $L{=}5$: reinforce $0.949$, TD $0.876$, no-trace $0.334$. $L{=}10$: $0.799$ / $0.793$ / $0.121$. $L{=}15$: $0.699$ / $0.641$ / $0.040$. $L{=}20$: $0.362$ / $0.224$ / $0.014$. \\
    \bottomrule
    \end{tabularx}
    \label{table:E10}
    \end{center}
    \vspace*{-\baselineskip}
\end{table}

% =====================================================================
\section{Device-grounded temporal networks (Chapter~\ref{ch:model})}

\noindent The following tables give the source data for the swept figures of the
device-grounded temporal networks and the working-memory experiments of
Chapter~\ref{ch:model} (the models reported in the \abbrev{APL} manuscript). The
capability and retention grids are replayed from
\texttt{data/results/temporal/*/temporal\_sweep.npy} and the store--recall grids from
\texttt{data/results/storerecall/*/storerecall.npy} of the \texttt{mnn-torch} package
(Appendix~\ref{app:code}); accuracies are percentages, delays and retention times in
abstract time-steps.

\begin{table}[H]
    \caption[Source data: benchmark capability (Fig.~\ref{fig:6cap}).]{Final test
    accuracy (\%) of the device-grounded temporal networks on the four streaming
    benchmarks, at the reference retention ($\tau_{leak}{=}30$, zero delay); the data of
    Figure~\ref{fig:6cap}. Mean\,[lo,\,hi] over $20$ seeds. Chance differs per task
    (column).}
    \begin{center}
    \footnotesize
    \begin{tabularx}{\textwidth}{@{}l *{4}{Y}@{}}
    \toprule
     & seq-MNIST & N-MNIST & SHD & DVS Gesture \\
     & (dense) & (conv) & (dense) & (conv) \\
    \midrule
    chance (\%)       & 10.0 & 10.0 & 5.0 & 9.1 \\
    accuracy (\%)     & 85.16 & 93.40 & 64.91 & 61.18 \\
    \;[lo,\,hi] & \tiny[84.84,85.45] & \tiny[93.22,93.59] & \tiny[64.09,65.76] & \tiny[60.02,62.22] \\
    \bottomrule
    \end{tabularx}
    \label{table:E11}
    \end{center}
    \vspace*{-\baselineskip}
\end{table}

\begin{table}[H]
    \caption[Source data: dense retention curves (Fig.~\ref{fig:6tr}a).]{Test accuracy
    (\%) of the dense streamed network (seq-MNIST) against appended delay $D$ (blank
    steps) for each membrane retention $\tau_{leak}$ and the no-memory control; the
    dense-panel data of Figure~\ref{fig:6tr}. Mean over $20$ seeds. Chance is $10\%$; the
    learned criterion is $44.96\%$.}
    \begin{center}
    \footnotesize
    \begin{tabularx}{\textwidth}{@{}l *{6}{Y}@{}}
    \toprule
    $D$ (steps) & 0 & 10 & 20 & 30 & 45 & 60 \\
    \midrule
    $\tau_{leak}{=}1.3$ & 49.8 & 21.4 & 13.1 & 13.3 & 12.1 & 12.3 \\
    $\tau_{leak}{=}2$   & 62.3 & 41.2 & 18.6 & 13.1 & 11.2 & 11.3 \\
    $\tau_{leak}{=}3$   & 69.5 & 53.0 & 33.9 & 18.5 & 15.4 & 11.2 \\
    $\tau_{leak}{=}5$   & 73.2 & 68.6 & 49.1 & 35.5 & 22.5 & 13.4 \\
    $\tau_{leak}{=}8$   & 74.0 & 73.4 & 60.3 & 45.9 & 32.7 & 24.3 \\
    $\tau_{leak}{=}12$  & 75.4 & 75.1 & 64.0 & 49.8 & 36.6 & 30.9 \\
    $\tau_{leak}{=}20$  & 77.6 & 75.2 & 67.2 & 57.2 & 41.8 & 34.5 \\
    $\tau_{leak}{=}30$  & 78.4 & 74.2 & 68.3 & 60.5 & 50.0 & 35.4 \\
    no memory           & 23.8 & 16.4 & 16.4 & 16.4 & 12.2 & 16.4 \\
    \bottomrule
    \end{tabularx}
    \label{table:E12}
    \end{center}
    \vspace*{-\baselineskip}
\end{table}

\begin{table}[H]
    \caption[Source data: convolutional retention curves (Fig.~\ref{fig:6tr}b).]{Test
    accuracy (\%) of the convolutional streamed network (N-MNIST) against appended delay
    $D$ (blank steps) for each membrane retention $\tau_{leak}$ and the no-memory
    control; the conv-panel data of Figure~\ref{fig:6tr}. Mean over $20$ seeds. Chance is
    $10\%$; the learned criterion is $50.38\%$.}
    \begin{center}
    \footnotesize
    \begin{tabularx}{\textwidth}{@{}l *{9}{Y}@{}}
    \toprule
    $D$ (steps) & 0 & 5 & 15 & 30 & 45 & 60 & 90 & 120 & 160 \\
    \midrule
    $\tau_{leak}{=}1.3$ & 87.0 & 81.8 & 21.2 & 10.0 & 10.0 & 10.0 & 10.0 & 10.0 & 10.0 \\
    $\tau_{leak}{=}2$   & 87.6 & 85.7 & 57.9 & 12.5 & 10.0 & 10.0 & 10.0 & 10.0 & 10.0 \\
    $\tau_{leak}{=}3$   & 88.7 & 87.0 & 72.2 & 33.9 & 13.8 & 10.0 & 10.0 & 10.0 & 10.0 \\
    $\tau_{leak}{=}5$   & 89.1 & 87.6 & 79.6 & 63.2 & 45.1 & 20.6 & 11.0 & 10.0 & 10.3 \\
    $\tau_{leak}{=}8$   & 89.6 & 87.7 & 83.0 & 71.8 & 63.8 & 56.2 & 20.5 & 13.2 & 10.2 \\
    $\tau_{leak}{=}12$  & 89.7 & 88.2 & 84.5 & 75.9 & 69.2 & 63.7 & 55.3 & 31.8 & 12.9 \\
    $\tau_{leak}{=}20$  & 89.8 & 88.1 & 85.3 & 79.0 & 73.0 & 68.7 & 62.1 & 56.4 & 40.3 \\
    $\tau_{leak}{=}30$  & 89.9 & 88.5 & 85.0 & 79.9 & 73.9 & 70.4 & 64.3 & 58.6 & 51.7 \\
    no memory           & 76.5 & 10.0 & 10.0 & 10.0 & 10.0 & 10.0 & 10.0 & 10.0 & 10.0 \\
    \bottomrule
    \end{tabularx}
    \label{table:E13}
    \end{center}
    \vspace*{-\baselineskip}
\end{table}

\begin{table}[H]
    \caption[Source data: store--recall working memory (Fig.~\ref{fig:6sr}).]{Recall
    accuracy (\%) in the store--recall (\abbrev{LSNN}) working-memory task against
    store--recall delay $D$ (blank steps), for the device membrane and the engineered
    \abbrev{ALIF} neuron at three retention timescales, with the fast-LIF (no-memory)
    control; the data of Figure~\ref{fig:6sr}. Three synapse rungs: (a) ideal, (b)
    healthy device Poole--Frenkel, (c) device under the measured stuck-off fault prior
    ($p{=}0.3$). Mean over $20$ seeds. Chance is $50\%$; the criterion is $75\%$.}
    \begin{center}
    \footnotesize
    \begin{tabularx}{\textwidth}{@{}l *{6}{Y}@{}}
    \toprule
    $D$ (steps) & 2 & 5 & 10 & 20 & 40 & 60 \\
    \midrule
    \multicolumn{7}{@{}l}{\emph{(a) ideal synapse}}\\
    fast LIF (no memory) & 100.0 & 100.0 & 99.8 & 57.9 & 52.4 & 53.0 \\
    device $\tau_{leak}{=}1.3$ & 100.0 & 100.0 & 49.9 & 49.9 & 50.1 & 49.8 \\
    device $\tau_{leak}{=}8$   & 100.0 & 100.0 & 100.0 & 100.0 & 90.0 & 94.0 \\
    device $\tau_{leak}{=}20$  & 100.0 & 100.0 & 100.0 & 100.0 & 99.9 & 99.0 \\
    ALIF $\tau_a{=}1.3$ & 100.0 & 100.0 & 100.0 & 56.3 & 53.7 & 56.2 \\
    ALIF $\tau_a{=}8$   & 100.0 & 100.0 & 100.0 & 100.0 & 100.0 & 51.7 \\
    ALIF $\tau_a{=}20$  & 100.0 & 100.0 & 100.0 & 100.0 & 100.0 & 100.0 \\
    \addlinespace
    \multicolumn{7}{@{}l}{\emph{(b) device synapse, healthy}}\\
    fast LIF (no memory) & 72.3 & 79.2 & 63.9 & 49.9 & 49.9 & 49.9 \\
    device $\tau_{leak}{=}1.3$ & 75.2 & 63.4 & 49.8 & 49.9 & 49.9 & 49.9 \\
    device $\tau_{leak}{=}8$   & 74.9 & 80.0 & 90.2 & 97.3 & 82.4 & 64.9 \\
    device $\tau_{leak}{=}20$  & 82.4 & 82.3 & 89.9 & 94.1 & 91.9 & 91.9 \\
    ALIF $\tau_a{=}1.3$ & 82.6 & 76.2 & 57.5 & 49.9 & 49.9 & 49.9 \\
    ALIF $\tau_a{=}8$   & 82.6 & 87.4 & 69.4 & 77.2 & 53.8 & 49.9 \\
    ALIF $\tau_a{=}20$  & 82.6 & 85.6 & 77.5 & 77.5 & 77.0 & 77.2 \\
    \addlinespace
    \multicolumn{7}{@{}l}{\emph{(c) device synapse, faulty ($p{=}0.3$ stuck-off)}}\\
    fast LIF (no memory) & 95.0 & 92.4 & 71.6 & 55.3 & 54.9 & 56.8 \\
    device $\tau_{leak}{=}1.3$ & 94.3 & 71.8 & 52.5 & 52.5 & 52.6 & 49.9 \\
    device $\tau_{leak}{=}8$   & 98.2 & 100.0 & 100.0 & 96.9 & 77.7 & 86.8 \\
    device $\tau_{leak}{=}20$  & 100.0 & 95.1 & 97.3 & 100.0 & 96.4 & 93.7 \\
    ALIF $\tau_a{=}1.3$ & 97.2 & 87.4 & 70.4 & 53.3 & 57.5 & 57.4 \\
    ALIF $\tau_a{=}8$   & 100.0 & 95.0 & 89.3 & 94.3 & 56.4 & 51.2 \\
    ALIF $\tau_a{=}20$  & 97.5 & 95.0 & 89.7 & 97.5 & 95.0 & 89.9 \\
    \bottomrule
    \end{tabularx}
    \label{table:E14}
    \end{center}
    \vspace*{-\baselineskip}
\end{table}

\begin{table}[H]
    \caption[Source data: homeostatic benefit by architecture (Fig.~\ref{fig:6homeo}).]{Homeostatic
    benefit (percentage points, homeostasis on minus off) by architecture; the data of
    Figure~\ref{fig:6homeo}. Mean\,[lo,\,hi] over $20$ seeds, clean ($p{=}0$). The dense
    classifier benefits; the temporal-convolutional network, already at ceiling on
    sustained activity, does not.}
    \begin{center}
    \footnotesize
    \begin{tabularx}{\textwidth}{@{}l *{3}{Y}@{}}
    \toprule
     & MSNN (dense) & MSNN (dense) & TemporalMCSNN \\
     & MNIST & Fashion & moving-MNIST \\
    \midrule
    benefit (pp) & $+2.62$ & $+4.36$ & $-8.5$ \\
    \;[lo,\,hi] & \tiny[1.98,3.38] & \tiny[3.22,5.61] & \tiny[-10.1,-6.8] \\
    \bottomrule
    \end{tabularx}
    \label{table:E15}
    \end{center}
    \vspace*{-\baselineskip}
\end{table}

\FloatBarrier
